% page 558 du fac-simile (image p-557 du scan)
% bloc 162.6 x 246.3 mm ; marge gauche 8.0 mm
\pgc{-12.700mm}{0.000mm}{
\l{\cel{3}550.}
\l{}
\l{\sou{stamino}. (\sou{bot.}) Organo sexuala maskulo, che la vejetanti.}
\l{\cel{3}- EFIS.}
\l{}
\l{\sou{stampar}. (\sou{trans}.) Provizar ye marko quan onu imprimas sur}
\l{\cel{3}objekto por atestar lua autentikeso. - DEFR.}
\l{}
\l{\sou{stano}. (\sou{kemio}).Korpo simpla, metalo grizatra, poke sonora,}
\l{\cel{3}poke duktila, tre forjebla, maleebla, de qua onu facas}
\l{\cel{3}implementi koqueyala. -\cel{1}\sou{Simbolo kemiala}\cel{1}: Sn. - EFIS.}
\l{}
\l{\sou{stanco}. Quanto fixa de versi di qua la senco esas kompleta,}
\l{\cel{3}dispozita segun ula ordino, e qua, repetate periodoze,}
\l{\cel{3}formacas la dividuri di poemeto. - DEFIRS.}
\l{}
\l{\sou{standar}. (\sou{netrans.}) I. (\sou{aludante persono o kozo}). Esar tala}
\l{\cel{3}en ica od en ita temp-uno. - II. (\sou{aludante persono)}.}
\l{\cel{3}Esar tala, en la hierarkio sociala. - DE.}
\l{}
\l{\sou{standardo}.\cel{2}La insigno nacionala, di qua un ek la bordi ver-}
\l{\cel{3}tikala muntesas sive an stango (armeo terala), sive an}
\l{\cel{3}kordo (armeo navala). -\cel{2}DEFIRS.}
\l{}
\l{\sou{stango}. Peco de ligno, de metalo, rigida, longa e streta.}
\l{\cel{3}DI.}
\l{}
\l{\sou{stapelio}. (\sou{bot.)}\cel{3}Genero de \textquotedbl{}asklepidacei\textquotedbl{}, ek planti xero-}
\l{\cel{3}fila, kun folii kaktusatra e granda flori. - L.\cel{1}\sou{stapelia}.}
\l{\cel{3}EFI?}
\l{}
\l{\sou{staplo}. I. Butikego ube onu vendas, ye la nomo di la stato}
\l{\cel{3}o di kompanio komercala. - II. Portuo, urbo komercala}
\l{\cel{3}ube la vari depozesas varte di lia exportaceso pos-impor-}
\l{\cel{3}taca. - DE.}
\l{}
\l{\sou{starchar}. (\sou{trans.}) I. Apretar per amilo. - II. (\sou{metaf.})}
\l{\cel{3}\sou{Starchita}\cel{1}: di qua la su-teno, la posturo, en mondumo,}
\l{\cel{3}esas tam rigida kam kamizo-plastrono apretita per amilo. DE}
\l{}
\l{\sou{startar}.\cel{2}I. (\sou{trans.}) Igar funcioneskar. - II. (\sou{netrans.})}
\l{\cel{3}Funcioneskar. - DE.}
\l{}
\l{\sou{startero}. (\sou{sur lico di keval-kuri}). La persono qua emisas la}
\l{\cel{3}departo-dgnalo, per inklinar la standardo quan lu tenas}
\l{\cel{3}en sua manuo dextra. - DEFIRS.}
\l{}
\l{\sou{stato}. La naciono egardata kom korpo politikala. - DEFIRS.}
\l{}
\l{\sou{statiko}. (\sou{fiziko}). Parto di mekaniko qua traktas la kondi-}
\l{\cel{3}cioni pri la equilibro di la forci. - DEFIRS.}
\l{}
\l{\sou{statistiko}. Cienco di qua la skopo esas donar dokumenti exak-}
\l{\cel{3}ta pri la parti diversa di la vivo di populo, per amasi-}
\l{\cel{3}gar, klasifikar, kontar omna fakti ek ta o ca naturo (naski,}
\l{\cel{3}morti, mariaji, produkturi, komerco, konsumo, kondamni, e c.)}
\l{\cel{3}qui eventas en periodo determinita. - DEFIRS.}
}
